\section{Implementation in the BX3 Framework}
\label{sec:implementation}

The Bailout Protocol is not a safety module appended to an existing agent runtime.
It is embedded in the BX3 execution model as a first-class component of the
three-layer architecture, sharing the same provisioning infrastructure, the same
deterministic execution engine, and the same append-only storage substrate as the
authorization ledger and the forensic ledger. This section specifies how the protocol
integrates with each BX3 layer, how the Bounds Engine drives bail-out triggering,
how the Fact Layer enforces the immutable audit trail, how the forensic ledger is
constructed, and how the state machine is embedded as a Fact Layer extension. It
establishes the formal relationship between the Bailout Protocol and BX3's upstream
accountability guarantee.

% ─────────────────────────────────────────────────────────────────
\subsection{Integration with the Purpose Layer: The Human Accountability Anchor}
\label{sec:purpose-layer}

At node initialization, the Purpose Layer of every BX3 node $n$ is provisioned with
two immutable structural elements:

\begin{itemize}
\item $\alpha(n)$ — the immutable parent pointer, establishing the node's position
  in the accountability chain.
\item $h(n)$ — the immutable reference to the human accountability anchor, a named
  human principal who bears legal and operational accountability for the node's
  autonomous actions.
\end{itemize}

Both elements are established at provisioning and cannot be altered at runtime by
any machine actor. The Purpose Layer carries $h(n)$ as a first-class structural
reference — not a metadata field, not a configuration comment. The Fact Layer's
authorization ledger maps $h(n)$ to the specific action classes the node is
permitted to execute, and the Bailout Protocol's escalation chain terminates at
$h(n)$ exclusively.

When a node enters the \texttt{ESCALATING} state, the escalation algorithm queries
the Purpose Layer to retrieve $h(n)$ and the provisioned communication pathways
registered to that principal. The Fact Layer records this retrieval in the forensic
ledger, creating an auditable chain from the trigger condition to the specific human
anchor who received the alert. No machine actor in the parent chain may substitute
for $h(n)$: the Purpose Layer does not proxy the request to a machine principal. It
returns $h(n)$ directly, and the communication pathway selection is made against the
pathways registered to $h(n)$ specifically.

The human anchor serves as the terminal resolution authority. When $h(n)$ issues a
directive via \texttt{HUMAN\_ACKNOWLEDGED}, that directive is bound to the trigger
package and propagated back down the escalation chain to the originating node. The
directive carries $h(n)$'s RDL identity, a cryptographic signature over the directive
payload, and a timestamp with non-repudiable provenance. The Fact Layer records the
directive as the definitive resolution of the trigger. No machine actor may modify,
supersede, or withdraw a human directive once issued; if the directive cannot be
executed, the node transitions to \texttt{RESOLVED\_EXECUTION\_FAILURE} and the Fact
Layer logs the discrepancy for post-hoc review.

The Purpose Layer enforces three structural properties that make human anchors
capable of terminating escalation chains: legal authority (human directives carry
legal weight that machine actors cannot replicate), scope fidelity (the
\texttt{ScopeOfAuthority} field ensures a human is never asked to resolve exceptions
outside their domain), and non-repudiability (cryptographic signatures are embedded
in the trigger package's immutable record).

The Training Wheels Protocol operates as a complementary, orthogonal accountability
channel to the Bailout Protocol. Training Wheels intercepts outbound actions before
they execute, providing prospective accountability; the Bailout Protocol receives
exception states that could not be resolved within the node's bounded authority,
providing retrospective accountability. Together they cover the full surface of node
behavior: Training Wheels governs what the node is \emph{authorized} to do; the
Bailout Protocol governs what happens when the node encounters a condition it was
not provisioned to handle.

% ─────────────────────────────────────────────────────────────────
\subsection{Bounds Engine Triggering Bail-Out}
\label{sec:bounds-trigger}

The Bounds Engine is the reasoning component of the BX3 three-layer model. It
maintains the \texttt{BoundsTable}, evaluates proposed actions against all registered
constraints, and determines whether a given action can proceed or must be escalated.
The Bailout Protocol is triggered exclusively through the Bounds Engine's
\texttt{BailoutCondition} evaluation, which runs continuously at each sensing cycle.

The Bounds Engine evaluates
$\texttt{BailoutCondition} \equiv \bigvee_{i=1}^{5} C_i$ at every sensing cycle.
Each condition $C_i$ maps to a distinct failure mode:

\begin{trivlist}\itemindent2em
\item[\textbf{$C_1$ — ConstraintViolation}] A proposed action violates one or more
  constraints in the \texttt{BoundsTable} and the violation cannot be resolved by
  constraint relaxation within the node's authority class.

\item[\textbf{$C_2$ — NovelInput}] The input state $x$ falls outside the training
  distribution $D_{\mathrm{train}}$ by a margin exceeding the node's novelty threshold
  $\tau_{\mathrm{novelty}}$, and no verified response strategy exists in the
  \texttt{ResponseTable}.

\item[\textbf{$C_3$ — ResourceExhaustion}] A required resource (time, memory, API
  quota, external dependency) is depleted such that safe completion of the current
  task cannot be guaranteed.

\item[\textbf{$C_4$ — HumanDirectiveConflict}] Two or more directives from distinct
  Purpose Layer anchors express conflicting intent, and the node lacks the authority
  to adjudicate.

\item[\textbf{$C_5$ — ForensicLedgerFailure}] A write to the Forensic Ledger fails
  or does not return a confirmation within the designated confirmation window. This
  condition is checked at every write operation; failure triggers immediate
  escalation because the audit trail cannot be guaranteed.
\end{trivlist}

When any $C_i$ evaluates to \texttt{True}, the Bounds Engine transitions the state
machine to \texttt{BOUNDED} and begins the local resolution procedure. The trigger
package $\mathcal{T}$ assembled at this point contains the condition sub-type, the
node identifier, the full state snapshot $\mathbf{s}_9$, and the current confidence
score $\phi$. The package is immutable from the moment of assembly — no machine actor
may modify $\mathcal{T}$ after it enters \texttt{BAIL\_PENDING}.

During \texttt{BOUNDED}, the Bounds Engine searches the action space $A$ for any
action $a$ satisfying
$\texttt{isPermissible}(a) \land \texttt{isSafe}(a) \land \texttt{hasResource}(a)$.
If such an $a$ exists, the node resolves the exception locally and transitions to
\texttt{IDLE}. If no such $a$ exists — formally,
$\neg \exists \; a \in A : \texttt{isPermissible}(a) \land \texttt{isSafe}(a) \land \texttt{hasResource}(a)$
— the Bounds Engine transitions to \texttt{BAIL\_PENDING}, assembling $\mathcal{T}$
for upstream propagation. The local resolution procedure is time-bounded by
$T(\texttt{BOUNDED}) = 30$ seconds; if the timeout fires before resolution, the
state transitions automatically to \texttt{BAIL\_PENDING} regardless of whether a
resolution candidate has been found.

In addition to TC, the Bounds Engine may emit an explicit \texttt{bailout\_signal}
at any time regardless of confidence scores. This handles conditions — for example,
a logical contradiction between two inputs both within $R(n)$ — that the Bounds
Engine recognizes as unresolvable even though each input is individually authorized.
The explicit signal bypasses TC entirely and immediately initiates the transition
to \texttt{BAIL\_PENDING}.

The Bounds Engine exercises no discretion over the trigger condition. The
evaluation of $\texttt{BailoutCondition}$ is a read-only query against the
BoundsTable. The result of the query directly drives the state machine transition
without an intervening decision step. There is no machine-actor gate between trigger
evaluation and protocol activation.

% ─────────────────────────────────────────────────────────────────
\subsection{Fact Layer Enforcing Audit Trail}
\label{sec:fact-layer}

The Fact Layer is the enforcement and audit component of the BX3 Framework. It is
responsible for deterministic action execution, cryptographic sealing of all state
transitions, and immutable maintenance of the Forensic Ledger. The Bailout Protocol
is deeply integrated with the Fact Layer in three ways: state machine enforcement,
cryptographic integrity, and tamper-evident chain construction.

The Bailout Protocol state machine $\mathcal{S}_{\text{bailout}}$ is not a separate
subsystem — it is an extension of the Fact Layer's enforcement authority. The Fact
Layer owns the state transition function $\delta$; no Bounds Engine or Purpose Layer
component may invoke a state transition directly. When the Bounds Engine evaluates a
trigger condition and determines that escalation is required, it \emph{requests} a
state transition from the Fact Layer. The Fact Layer validates the transition request
against the current state and event, executes the transition atomically, and seals
the result in the Forensic Ledger before acknowledging the transition to the requesting
component.

This enforcement pattern ensures that the state machine cannot be manipulated by any
component outside the Fact Layer. A machine actor that attempts to skip states, forge
a transition, or fabricate a trigger condition will fail at the Fact Layer's validation
step, producing a $C_5$ (\texttt{ForensicLedgerFailure}) condition that triggers its
own escalation.

Every state transition in $\mathcal{S}_{\text{bailout}}$ produces a Forensic Ledger
entry with the following structure:

\begin{equation}
\mathcal{L}_{\text{bailout}} = \bigl\langle
  \tau_{\text{id}},         % Trigger ID (globally unique)
  q_{\text{from}},           % Source state
  q_{\text{to}},             % Target state
  e,                         % Triggering event
  n_{\text{actor}},          % Node executing transition
  t_{\text{transition}},     % Wall-clock timestamp (nanosecond precision)
  \mathbf{s}_9,            % State snapshot at transition
  \phi,                      % Confidence score at fire time
  \text{digest},            % SHA-256: H(this_entry ∥ previous_digest)
  \rho_{\text{final}}        % Final resolution (populated at RESOLVED)
\bigr\rangle
\end{equation}

The digest field implements a chain integrity mechanism. Each entry's digest is
computed as the SHA-256 hash of the concatenation of the entry's raw fields and the
previous entry's digest. This creates a tamper-evident chain: modifying any
historical entry breaks the digest chain from that point forward, and the break is
detectable by any independent auditor without access to the chain's internal state.

The Fact Layer enforces the Audit Completeness Guarantee through a mandatory write
confirmation requirement. A state transition is not considered complete until the
Forensic Ledger has confirmed the write with a positive acknowledgment. If the write
fails or the acknowledgment is not received within the confirmation window, the Fact
Layer sets $C_5$ and triggers an autonomous escalation from the current state. The
node does not proceed with any action that depends on an unconfirmed state
transition.

% ─────────────────────────────────────────────────────────────────
\subsection{Forensic Ledger Recording}
\label{sec:forensic-ledger}

The forensic ledger is the append-only, cryptographically sealed record of every
protocol event in the Bailout Protocol state machine. It is maintained by the Fact
Layer as a first-class data structure — not a debug log, not a side channel, not an
optional diagnostic output. Every state transition, trigger evaluation, timeout
event, pathway selection, human directive, and resolution action is written to the
ledger as an immutable entry. The ledger is the authoritative source of truth for
protocol compliance verification.

Every Bailout Protocol event produces one or more Forensic Ledger entries across
three planes, corresponding to the BX3 Framework's three-layer architecture:

\begin{trivlist}\itemindent2em
\item[\textbf{Purpose Plane}] Records authorization to escalate, human acknowledgment
  events, directive issuance, and principal identity at each escalation hop. These
  entries establish who knew about the exception and what they directed.

\item[\textbf{Bounds Engine Plane}] Records trigger condition evaluations ($C_1$
  through $C_5$), local resolution attempts, constraint violation details,
  confidence scores at fire time, and escalation path metadata. These entries
  establish what the machine observed and why it concluded escalation was necessary.

\item[\textbf{Fact Plane}] Records state transition timestamps, timeout fire events,
  transition validation confirmations, resolution action executions, and Forensic
  Ledger write acknowledgments. These entries establish what the system did at each
  step and that the audit trail itself is intact.
\end{trivlist}

The Forensic Ledger entries for a complete trigger lifecycle form a directed acyclic
graph rooted at the originating node's first \texttt{BOUNDED} entry. The graph
captures the full provenance of the exception: initial condition, local resolution
attempt, escalation path (all hops), human acknowledgment(s), directive (if issued),
actuator execution (or execution failure), and closure. No entry in this graph can
be deleted, modified, or reordered after commitment. The chain is append-only and
cryptographically sealed via the SHA-256 digest chain described above.

The forensic ledger is what distinguishes the Bailout Protocol from conventional
exception handling. A conventional error-handling routine might log an error, might
notify an operator, and might attempt to recover. But none of these actions are
architecturally compelled, none are immutably recorded, and none guarantee human
receipt. The Bailout Protocol's forensic ledger converts the protocol's guarantees
from operational assertions into independently verifiable evidence. The system did
or did not escalate; the human anchor did or did not receive the exception; the
human principal did or did not issue a directive. These facts are recorded, sealed,
and attributable. The ledger is the accountability artifact.

% ─────────────────────────────────────────────────────────────────
\subsection{State Machine as Fact Layer Extension}
\label{sec:state-machine}

The deterministic state machine $\mathcal{B}(n)$ is embedded in the Fact Layer as
an extension of its enforcement function. The Fact Layer's core property is
deterministic enforcement: it enforces constraints without judgment, without
discretion, and without the ability to be bypassed by any machine actor. The state
machine extends this property from static constraint enforcement to dynamic protocol
enforcement: it governs not merely whether an action may execute, but in what
sequence and under what conditions actions may execute, and what the mandatory
outcome is when conditions exceed the node's provisioned scope.

The state $q \in Q$ of a node's Bailout Protocol machine is stored in the node's
Fact Layer worksheet — the same protected, atomic storage used for the authorization
ledger and the context package. State transitions are executed by the Fact Layer's
pre-commit hook, which evaluates the transition relation before committing any state
change. The pre-commit hook runs synchronously before the node's next scheduling
quantum is allocated, ensuring that no state transition can be deferred, suppressed,
or altered by the node's own execution thread.

The transition relation $\rightarrow \subseteq Q \times \Sigma \times Q$ is defined
as a lookup table in the Fact Layer's deterministic execution engine. The priority
function $\pi(n, x)$ resolves simultaneous trigger conditions by totally ordering them
before transition commitment:

\begin{equation}
\pi(n, x) \;=\; w_1 \cdot \text{severity}(x) \;+\; w_2 \cdot \bigl(1 - \text{confidence}(n, x)\bigr) \;+\; w_3 \cdot \text{distance}\bigl(n, h(n)\bigr)
\tag{PF}
\end{equation}

where $w_1 > w_2 > w_3 \geq 0$ are provisioned constants. The priority function is
evaluated by the Bailout Monitor — a Fact Layer component — before any state change
is committed. No machine actor can override $\pi$.

The state machine's deterministic property is enforced at the architecture level:
for every $(q, \sigma) \in Q \times \Sigma$, at most one $q' \in Q$ satisfies
$(q, \sigma, q') \in \rightarrow$. If no transition is defined for a $(q, \sigma)$
pair, the machine stalls and the Bailout Monitor records a protocol violation in
the forensic ledger. Stall conditions are terminal audit events indicating a
protocol implementation error or a state corruption that requires human remediation.

Embedding $\mathcal{B}(n)$ in the Fact Layer provides four guarantees that a
standalone or middleware implementation cannot:

\begin{enumerate}
\item \textbf{State isolation:} The state $q$ is stored in the Fact Layer worksheet
  and is observable only by the Fact Layer and its designated monitors. No machine
  actor operating at user level or elevated privilege can read or write $q$ directly.
\item \textbf{Transition atomicity:} State transitions are atomic from the
  scheduler's perspective. No intermediate or partially-executing states are
  observable externally. The transition either completes in full or does not occur
  at all.
\item \textbf{Deterministic evaluation:} The transition relation and the priority
  function are executed by the Fact Layer's deterministic engine. Their outputs are
  determined solely by the current state and the input event; no stochastic or
  heuristic decision-making influences the outcome.
\item \textbf{Append-only audit:} All transition events are written to the forensic
  ledger immediately before the state change is committed. The ledger entry is
  created atomically with the state update, ensuring that no transition occurs
  without a corresponding audit record.
\end{enumerate}

% ─────────────────────────────────────────────────────────────────
\subsection{Bailout Protocol as Runtime Expression of the Upstream Accountability Guarantee}
\label{sec:runtime-expression}

The BX3 Framework makes a specific architectural claim: that autonomous systems
built on its three-layer architecture enforce an \emph{upstream accountability
guarantee}. The guarantee states that any BX3 node encountering a state it cannot
resolve within its provisioned operating range must escalate accountability
recursively upward through the system hierarchy until it reaches a human principal
— and that no machine actor along that escalation chain may absorb, redirect, or
suppress the exception. This guarantee is a design-time property of the framework.

The Bailout Protocol is the runtime expression of this guarantee. The guarantee
describes what BX3 systems are \emph{obligated} to do; the Bailout Protocol is the
mechanism that \emph{compels} them to do it. The distinction is between a
specification and an implementation: the guarantee is the specification of the
accountability property, and the Bailout Protocol is the implementation that makes
the specification enforceable.

This relationship is structurally necessary, not accidentally related. The upstream
accountability guarantee requires three conditions to hold simultaneously:
(1) a human principal must be designated as the terminus of every escalation chain;
(2) the escalation path must contain no machine actor as a terminus;
(3) every step of the escalation must be immutably recorded.
The Bailout Protocol implements all three conditions:
condition (1) is implemented by the immutable $h(n)$ reference set at node
provisioning;
condition (2) is implemented by the Fact Layer's pre-commit hook and the immutable
parent-pointer architecture;
condition (3) is implemented by the forensic ledger maintained by the Fact Layer.

Removing any layer from the BX3 three-layer model renders the Bailout Protocol
incomplete. Without the Purpose Layer, there is no accountability terminus — no
human principal to receive escalations. Without the Bounds Engine, there is no
authoritative source of scope judgments — no principled way to determine what falls
inside or outside the node's provisioned operating range. Without the Fact Layer,
there is no enforcement mechanism for the protocol's compulsions and no immutable
recording medium for the forensic ledger. The three-layer architecture is not merely
convenient; it is structurally necessary for the upstream accountability guarantee.

This necessary relationship is what makes the Bailout Protocol architecturally
distinct from conventional escalation hierarchies. Most AI systems implement some
form of human escalation for critical errors. The Bailout Protocol is different in
being structurally compelled: the escalation is not a fallback that the system
chooses when things go wrong, but the prescribed resolution path for any
unresolvable state at any depth. The parent pointer chain terminates at a human not
because the system decided to escalate, but because the architecture defines
escalation as the only permissible outcome for conditions that exceed the Bounds
Engine's provisioned scope.

The upstream accountability guarantee and the Bailout Protocol are two aspects of
the same architectural commitment. The guarantee is the promise; the Bailout Protocol
is the mechanism that fulfills that promise. Together they establish the core safety
property of the BX3 Framework: that autonomous capability in BX3 systems is always
bounded by human accountability, and that the bounding is enforced architecturally
rather than operationally. The system fails upward into human consciousness — never
downward into autonomous chaos.

% ─────────────────────────────────────────────────────────────────
\subsection{Summary of Formal Properties}
\label{sec:properties-summary}

The implementation of the Bailout Protocol within the BX3 Framework satisfies the
following formal properties:

\begin{trivlist}\itemindent0em
\item[\textbf{Upstream Accountability Guarantee}] For any node $n$, there exists a
  finite escalation path $P = \langle n = p_0, p_1, \ldots, p_k \rangle$ such that
  $p_k$ is a Human Accountability Anchor, and $\forall i \in [0, k-1]$,
  $\texttt{IsMachineActor}(p_i) = \text{True}$. The path is constructed by the
  \texttt{GetNextUpstreamAnchor} algorithm operating on the Purpose Layer anchor
  registry.

\item[\textbf{Non-Orphaning Guarantee}] No trigger can reach \texttt{RESOLVED} without
  at least one human acknowledgment in its escalation path, unless the resolution is
  a timeout exception at \texttt{MaxEscalationDepth}, in which case the node enters
  \texttt{HALT\_REQUIRED} state and cannot resume autonomous operation without human
  clearance.

\item[\textbf{Deterministic State Transitions}] Every transition in
  $\mathcal{S}_{\text{bailout}}$ is deterministic and enforced by the Fact Layer.
  No machine actor may invoke, skip, or modify a state transition.

\item[\textbf{Audit Completeness Guarantee}] Every state transition produces a
  confirmed Forensic Ledger entry across all three planes (Purpose, Bounds Engine,
  Fact) before the transition is considered complete. A transition without a
  confirmed write is invalid and triggers $C_5$.

\item[\textbf{Immutable Trigger Package}] After the Bounds Engine assembles the
  trigger package $\mathcal{T}$ at \texttt{BOUNDED}, no machine actor may modify any
  field of $\mathcal{T}$. Modifications are detectable at Fact Layer validation and
  produce a forensic anomaly entry.

\item[\textbf{Graceful Degradation}] Network disconnection does not suppress or delay
  escalation. Triggers generated during disconnection are buffered at the originating
  node and delivered upon reconnection, preserving the full escalation path in the
  Forensic Ledger.
\end{trivlist}

These properties collectively ensure that the Bailout Protocol, as implemented within
the BX3 Framework, delivers the BX3 Framework's core safety promise in all operational
contexts: \emph{accountability always escalates to a human, or the system halts and
awaits human review.}